\documentclass[12pt,a4paper]{article}

% ------------------------------------------------
% Paquetes
% ------------------------------------------------
\usepackage[utf8]{inputenc}
\usepackage[T1]{fontenc}
\usepackage[spanish, es-nodecimaldot]{babel}
\usepackage{geometry}
\usepackage{graphicx}
\usepackage{float}
\usepackage{xcolor}
\usepackage{listings}
\usepackage{hyperref}
\usepackage{fancyhdr}
\usepackage{titlesec}
\usepackage{amsmath}

% ------------------------------------------------
% Configuración de página
% ------------------------------------------------
\geometry{
    top=2.5cm,
    bottom=2.5cm,
    left=2.5cm,
    right=2.5cm
}

% ------------------------------------------------
% Configuración de listings
% ------------------------------------------------
\definecolor{codegray}{gray}{0.95}

\lstset{
    backgroundcolor=\color{codegray},
    basicstyle=\ttfamily\small,
    breaklines=true,
    frame=single,
    columns=fullflexible,
    keepspaces=true,
    showstringspaces=false,
    tabsize=2
}

% ------------------------------------------------
% Configuración de títulos
% ------------------------------------------------
\titleformat{\section}
    {\large\bfseries}
    {\thesection.}
    {0.5em}
    {}

% ------------------------------------------------
% Hipervínculos
% ------------------------------------------------
\hypersetup{
    colorlinks=true,
    linkcolor=black,
    urlcolor=blue,
    citecolor=black
}

% ------------------------------------------------
% Comando para insertar capturas opcionales
% ------------------------------------------------
\newcommand{\captura}[2]{%
    \begin{figure}[H]
        \centering
        \IfFileExists{#1}{%
            \includegraphics[
                width=0.92\textwidth,
                keepaspectratio
            ]{#1}
        }{%
            \fbox{%
                \parbox[c][4cm][c]{0.82\textwidth}{%
                    \centering
                    \textbf{Captura pendiente}\\[0.5em]
                    Archivo esperado: \texttt{#1}
                }%
            }%
        }
        \caption{#2}
    \end{figure}
}

% ------------------------------------------------
% Datos del documento
% ------------------------------------------------
\title{
    MPI:\\[0.3em]
    \textbf{Suma de dos arreglos}\\[0.5em]
    \large Supercómputo
}

\author{
    Jorge Alberto Suárez Saldaña\\
    \small Matrícula: 355992\\
    \small Profesor: Raymundo Antonio González Grimaldo
}

\date{\today}

% ------------------------------------------------
% Documento
% ------------------------------------------------
\begin{document}

\maketitle

\section{Introducción}

En la presente tarea se utilizó la interfaz de paso de mensajes (MPI) en el
lenguaje de programación C para paralelizar la suma de dos arreglos de números
enteros. Se implementó un esquema coordinador--trabajadores en el que el proceso
con identificador cero genera los arreglos originales \texttt{a} y \texttt{b},
distribuye sus elementos entre los demás procesos y recibe las secciones del
arreglo resultante \texttt{c}.

Cada trabajador suma elemento a elemento las secciones que recibe y devuelve
sus resultados al proceso coordinador. Finalmente, este reconstruye el arreglo
completo y verifica que cada elemento cumpla la igualdad
\(c[i]=a[i]+b[i]\). Este ejercicio permite estudiar el reparto de datos, las
comunicaciones punto a punto y la coordinación de procesos en una aplicación
paralela.

\section{Código fuente}

El siguiente programa se escribió en C utilizando las funciones de MPI para
el envío y la recepción de mensajes.

\begin{lstlisting}[language=C]
/* Materia: Supercomputo
 * Semestre: Noveno
 * Nombre del alumno: Jorge Alberto Suarez Saldana
 * Clave del alumno: 355992
 * Carrera: Ingenieria en Sistemas Inteligentes
 * Nombre de la tarea: Suma de dos arreglos
 * Periodo de evaluacion: 2do parcial
 * Descripcion: Suma los elementos de los arreglos a y b y
 *              los guarda en el arreglo c utilizando MPI
 * Avance logrado: 100%
 */

#include <stdio.h>
#include <stdlib.h>
#include <time.h>
#include <mpi.h>

#define N 100

int main(int argc, char **argv) {
    int idProc, numProc;
    int nDatos, nDatosU;
    int *a, *b, *c;
    int inicio;

    MPI_Init(&argc, &argv);

    MPI_Comm_rank(MPI_COMM_WORLD, &idProc);
    MPI_Comm_size(MPI_COMM_WORLD, &numProc);

    if (numProc < 2 || numProc > N + 1) {
        if (idProc == 0)
            printf("Se requieren entre 2 y %d procesos.\n", N + 1);

        MPI_Finalize();
        return 1;
    }

    nDatos = N / (numProc - 1);
    nDatosU = nDatos + N % (numProc - 1);

    if (idProc == 0) {
        // Proceso coordinador
        a = malloc(N * sizeof(int));
        b = malloc(N * sizeof(int));
        c = malloc(N * sizeof(int));

        if (a == NULL || b == NULL || c == NULL) {
            MPI_Abort(MPI_COMM_WORLD, 1);
        }

        // Llenar arreglos con numeros aleatorios
        srand(time(NULL));

        for (int i = 0; i < N; i++) {
            a[i] = rand() % 100 + 1;
            b[i] = rand() % 100 + 1;
        }

        // Enviar secciones a los trabajadores
        for (int iProc = 1; iProc < numProc - 1; iProc++) {
            inicio = (iProc - 1) * nDatos;

            MPI_Send(&a[inicio], nDatos, MPI_INT,
                     iProc, 0, MPI_COMM_WORLD);

            MPI_Send(&b[inicio], nDatos, MPI_INT,
                     iProc, 1, MPI_COMM_WORLD);
        }

        // Enviar ultima seccion
        inicio = (numProc - 2) * nDatos;

        MPI_Send(&a[inicio], nDatosU, MPI_INT,
                 numProc - 1, 0, MPI_COMM_WORLD);

        MPI_Send(&b[inicio], nDatosU, MPI_INT,
                 numProc - 1, 1, MPI_COMM_WORLD);

        // Recibir resultados de los trabajadores
        for (int iProc = 1; iProc < numProc - 1; iProc++) {
            inicio = (iProc - 1) * nDatos;

            MPI_Recv(&c[inicio], nDatos, MPI_INT,
                     iProc, 2, MPI_COMM_WORLD,
                     MPI_STATUS_IGNORE);
        }

        // Recibir ultima seccion de c
        inicio = (numProc - 2) * nDatos;

        MPI_Recv(&c[inicio], nDatosU, MPI_INT,
                 numProc - 1, 2, MPI_COMM_WORLD,
                 MPI_STATUS_IGNORE);

        // Imprimir los tres arreglos
        printf("Arreglo A:\n");
        for (int i = 0; i < N; i++) {
            printf("%d ", a[i]);
        }

        printf("\n\nArreglo B:\n");
        for (int i = 0; i < N; i++) {
            printf("%d ", b[i]);
        }

        printf("\n\nArreglo C (A + B):\n");
        for (int i = 0; i < N; i++) {
            printf("%d ", c[i]);
        }

        printf("\n\n");

        // Verificar el resultado
        int correcto = 1;

        for (int i = 0; i < N; i++) {
            if (c[i] != a[i] + b[i]) {
                correcto = 0;
                break;
            }
        }

        if (correcto)
            printf("Resultado correcto.\n");
        else
            printf("Error: los resultados no coinciden.\n");

        free(a);
        free(b);
        free(c);

    } else {
        // Procesos trabajadores
        if (idProc == numProc - 1) {
            nDatos = nDatosU;
        }

        a = malloc(nDatos * sizeof(int));
        b = malloc(nDatos * sizeof(int));
        c = malloc(nDatos * sizeof(int));

        if (a == NULL || b == NULL || c == NULL) {
            MPI_Abort(MPI_COMM_WORLD, 1);
        }

        // Recibir secciones de a y b
        MPI_Recv(a, nDatos, MPI_INT,
                 0, 0, MPI_COMM_WORLD,
                 MPI_STATUS_IGNORE);

        MPI_Recv(b, nDatos, MPI_INT,
                 0, 1, MPI_COMM_WORLD,
                 MPI_STATUS_IGNORE);

        // Sumar los arreglos elemento a elemento
        for (int i = 0; i < nDatos; i++) {
            c[i] = a[i] + b[i];
        }

        // Devolver resultado al coordinador
        MPI_Send(c, nDatos, MPI_INT,
                 0, 2, MPI_COMM_WORLD);

        free(a);
        free(b);
        free(c);
    }

    MPI_Finalize();
    return 0;
}
\end{lstlisting}

\section{Explicación del código}

El programa incluye las bibliotecas \texttt{stdio.h} para entrada y salida,
\texttt{stdlib.h} para memoria dinámica y generación de números aleatorios,
\texttt{time.h} para obtener una semilla y \texttt{mpi.h} para acceder a las
funciones de comunicación entre procesos.

La constante \texttt{N} determina la cantidad de elementos de cada arreglo:

\begin{lstlisting}[language=C]
#define N 100

MPI_Init(&argc, &argv);
MPI_Comm_rank(MPI_COMM_WORLD, &idProc);
MPI_Comm_size(MPI_COMM_WORLD, &numProc);
\end{lstlisting}

\texttt{MPI\_Init()} inicializa el entorno MPI. La función
\texttt{MPI\_Comm\_rank()} obtiene el identificador del proceso actual y lo
almacena en \texttt{idProc}, mientras que \texttt{MPI\_Comm\_size()} almacena
el número total de procesos en \texttt{numProc}. El comunicador
\texttt{MPI\_COMM\_WORLD} incluye inicialmente a todos los procesos de la
ejecución.

El programa exige al menos dos procesos, pues necesita un coordinador y un
trabajador. También limita la ejecución a \(N+1=101\) procesos para evitar
trabajadores sin elementos que procesar.

\subsection{Distribución de los datos}

El proceso cero se reserva para coordinar la ejecución, mientras que los
\(P-1\) procesos restantes se encargan de realizar las sumas. La cantidad de
elementos que recibe cada trabajador, excepto el último, se calcula mediante:

\begin{lstlisting}[language=C]
nDatos = N / (numProc - 1);
nDatosU = nDatos + N % (numProc - 1);
\end{lstlisting}

Si \(P\) es el número total de procesos, el número de trabajadores es:

\[
P_t=P-1.
\]

Para los primeros \(P-2\) trabajadores, el tamaño de la sección es:

\[
nDatos=\left\lfloor\frac{N}{P-1}\right\rfloor.
\]

El último trabajador procesa la misma cantidad más el residuo de la división:

\[
nDatosU=\left\lfloor\frac{N}{P-1}\right\rfloor
+\bigl(N\bmod(P-1)\bigr).
\]

Por ejemplo, al ejecutar el programa con siete procesos, seis de ellos serán
trabajadores. Para \(N=100\):

\[
nDatos=\left\lfloor\frac{100}{6}\right\rfloor=16,
\qquad
nDatosU=16+(100\bmod 6)=20.
\]

Por lo tanto, cinco trabajadores recibirán 16 elementos de cada arreglo y el
último recibirá 20. Así se procesarán los \(5(16)+20=100\) elementos de cada
arreglo sin perder información.

\subsection{Funcionamiento del proceso coordinador}

El proceso con identificador cero reserva memoria para tres arreglos de
\(N\) enteros: \texttt{a}, \texttt{b} y \texttt{c}. Los dos primeros contienen
los operandos, mientras que el tercero almacenará la suma.

\begin{lstlisting}[language=C]
a = malloc(N * sizeof(int));
b = malloc(N * sizeof(int));
c = malloc(N * sizeof(int));

if (a == NULL || b == NULL || c == NULL) {
    MPI_Abort(MPI_COMM_WORLD, 1);
}
\end{lstlisting}

Si alguna reserva falla, \texttt{MPI\_Abort()} solicita la interrupción de la
ejecución de los procesos pertenecientes al comunicador.

\subsubsection{Generación de los arreglos originales}

Únicamente el coordinador llena los arreglos \texttt{a} y \texttt{b} con
valores pseudoaleatorios enteros entre 1 y 100:

\begin{lstlisting}[language=C]
srand(time(NULL));

for (int i = 0; i < N; i++) {
    a[i] = rand() % 100 + 1;
    b[i] = rand() % 100 + 1;
}
\end{lstlisting}

La instrucción \texttt{srand()} inicializa la semilla del generador usando la
hora actual. Los trabajadores no generan números aleatorios: solamente
procesan los datos que reciben.

\subsubsection{Envío de las secciones de A y B}

El coordinador calcula la posición inicial de cada segmento y utiliza dos
llamadas a \texttt{MPI\_Send()} para enviar las secciones correspondientes de
los dos arreglos:

\begin{lstlisting}[language=C]
for (int iProc = 1; iProc < numProc - 1; iProc++) {
    inicio = (iProc - 1) * nDatos;

    MPI_Send(&a[inicio], nDatos, MPI_INT,
             iProc, 0, MPI_COMM_WORLD);

    MPI_Send(&b[inicio], nDatos, MPI_INT,
             iProc, 1, MPI_COMM_WORLD);
}
\end{lstlisting}

La expresión
\[
inicio=(iProc-1)\cdot nDatos
\]
identifica el índice a partir del cual se encuentran los datos que recibirá
cada trabajador. Las secciones de \texttt{a} y \texttt{b} comienzan en el
mismo índice, pues se sumarán elemento a elemento.

Cada llamada a \texttt{MPI\_Send()} especifica la dirección del primer
elemento, la cantidad de enteros, el tipo \texttt{MPI\_INT}, el identificador
del destinatario, una etiqueta y el comunicador. En este programa se utilizan
las etiquetas \texttt{0} para las secciones de \texttt{a} y \texttt{1} para
las secciones de \texttt{b}.

El último trabajador requiere un envío separado, porque puede recibir más
elementos que los anteriores:

\begin{lstlisting}[language=C]
inicio = (numProc - 2) * nDatos;

MPI_Send(&a[inicio], nDatosU, MPI_INT,
         numProc - 1, 0, MPI_COMM_WORLD);

MPI_Send(&b[inicio], nDatosU, MPI_INT,
         numProc - 1, 1, MPI_COMM_WORLD);
\end{lstlisting}

De esta manera, cada trabajador recibe exclusivamente las secciones que debe
procesar, y no los arreglos completos.

\subsubsection{Recepción y reconstrucción del arreglo C}

Después del envío, el coordinador recibe de cada trabajador la sección
resultante de \texttt{c}, utilizando la etiqueta \texttt{2}:

\begin{lstlisting}[language=C]
for (int iProc = 1; iProc < numProc - 1; iProc++) {
    inicio = (iProc - 1) * nDatos;

    MPI_Recv(&c[inicio], nDatos, MPI_INT,
             iProc, 2, MPI_COMM_WORLD,
             MPI_STATUS_IGNORE);
}

inicio = (numProc - 2) * nDatos;

MPI_Recv(&c[inicio], nDatosU, MPI_INT,
         numProc - 1, 2, MPI_COMM_WORLD,
         MPI_STATUS_IGNORE);
\end{lstlisting}

A diferencia de un algoritmo que solo devuelve una suma total, aquí cada
trabajador envía un subarreglo de resultados. El coordinador almacena cada
mensaje directamente en la posición correspondiente del arreglo \texttt{c}.
Por ello, las secciones llegan a ocupar índices consecutivos y se reconstruye
el arreglo completo sin operaciones adicionales de copia.

\texttt{MPI\_Recv()} recibe un mensaje del proceso y con la etiqueta
especificados. La constante \texttt{MPI\_STATUS\_IGNORE} indica que no se
necesita guardar los datos de estado de la recepción.

\subsection{Funcionamiento de los procesos trabajadores}

Todos los procesos distintos de cero ejecutan el bloque destinado a los
trabajadores. Si el proceso es el último, ajusta la cantidad de elementos que
debe recibir:

\begin{lstlisting}[language=C]
if (idProc == numProc - 1) {
    nDatos = nDatosU;
}
\end{lstlisting}

Después reserva memoria independiente para sus tres arreglos locales y
recibe las secciones enviadas por el coordinador:

\begin{lstlisting}[language=C]
MPI_Recv(a, nDatos, MPI_INT,
         0, 0, MPI_COMM_WORLD,
         MPI_STATUS_IGNORE);

MPI_Recv(b, nDatos, MPI_INT,
         0, 1, MPI_COMM_WORLD,
         MPI_STATUS_IGNORE);
\end{lstlisting}

Ambos mensajes provienen del proceso cero. Las etiquetas \texttt{0} y
\texttt{1} permiten distinguir los datos del primer y del segundo arreglo.
Cada proceso dispone de su propio espacio de memoria, por lo que sus arreglos
locales son independientes de los del coordinador.

\subsubsection{Suma elemento a elemento}

Una vez recibidos ambos segmentos, el trabajador calcula cada elemento de su
arreglo \texttt{c} mediante el siguiente ciclo:

\begin{lstlisting}[language=C]
for (int i = 0; i < nDatos; i++) {
    c[i] = a[i] + b[i];
}
\end{lstlisting}

La operación realizada para cada posición se expresa como:

\[
c_i=a_i+b_i,\qquad 0\leq i<N.
\]

Aunque cada trabajador recorre secuencialmente su propia sección, los
diferentes procesos pueden realizar los cálculos concurrentemente. Esto
constituye una paralelización de datos: la misma operación se aplica a
subconjuntos distintos de los arreglos originales.

\subsubsection{Devolución de los resultados}

Finalmente, el trabajador envía al coordinador todos los valores de su
subarreglo \texttt{c}:

\begin{lstlisting}[language=C]
MPI_Send(c, nDatos, MPI_INT,
         0, 2, MPI_COMM_WORLD);
\end{lstlisting}

La etiqueta \texttt{2} distingue el mensaje de resultado de los mensajes de
entrada. Una vez realizado el envío, el trabajador libera sus tres bloques de
memoria con \texttt{free()}.

\subsection{Organización de los mensajes}

Para distinguir los tipos de datos intercambiados se emplean tres etiquetas.
Su función dentro del programa es la siguiente:

\begin{center}
\begin{tabular}{|c|l|l|}
\hline
\textbf{Etiqueta} & \textbf{Origen y destino} & \textbf{Contenido} \\
\hline
0 & Coordinador $\rightarrow$ trabajador & Sección de \texttt{a} \\
\hline
1 & Coordinador $\rightarrow$ trabajador & Sección de \texttt{b} \\
\hline
2 & Trabajador $\rightarrow$ coordinador & Sección de \texttt{c} \\
\hline
\end{tabular}
\end{center}

Cada trabajador recibe dos mensajes y envía uno. Por lo tanto, con \(P-1\)
trabajadores se realizan \(2(P-1)\) envíos de segmentos de entrada y
\(P-1\) envíos de resultados, es decir, \(3(P-1)\) mensajes de datos en total.
En conjunto, se transfieren \(2N\) elementos enteros hacia los trabajadores
y \(N\) elementos enteros de regreso al coordinador, sin contar información
interna del protocolo de comunicación.

\subsection{Impresión y verificación de resultados}

Cuando el proceso cero ya ha recibido todas las secciones, imprime los tres
arreglos completos, permitiendo examinar los datos de entrada y el resultado.
Después comprueba que cada elemento de \texttt{c} sea igual a la suma de los
valores correspondientes en \texttt{a} y \texttt{b}:

\begin{lstlisting}[language=C]
int correcto = 1;

for (int i = 0; i < N; i++) {
    if (c[i] != a[i] + b[i]) {
        correcto = 0;
        break;
    }
}

if (correcto)
    printf("Resultado correcto.\n");
else
    printf("Error: los resultados no coinciden.\n");
\end{lstlisting}

La variable \texttt{correcto} empieza con valor \texttt{1}, que representa
que no se han detectado errores. Si alguna posición no cumple la igualdad,
se asigna el valor \texttt{0} y termina la revisión mediante \texttt{break}.
Si todas las posiciones coinciden, se muestra el mensaje
\texttt{Resultado correcto.}

Esta validación permite detectar errores en el reparto, el cálculo o la
reconstrucción de los datos. Comprueba la corrección para los valores usados
en esa ejecución, aunque no sustituye una evaluación del rendimiento del
algoritmo paralelo.

\subsection{Finalización y liberación de recursos}

Tanto el coordinador como los trabajadores liberan con \texttt{free()} los
bloques de memoria reservados. Al concluir la comunicación, todos los
procesos ejecutan:

\begin{lstlisting}[language=C]
MPI_Finalize();
\end{lstlisting}

Esta función finaliza ordenadamente el entorno MPI y libera los recursos
internos asociados a la ejecución.

\section{Compilación y ejecución}

Para compilar el programa, guardado con el nombre
\texttt{sumaDosArreglos.c}, se utiliza el compilador de MPI:

\begin{lstlisting}
mpicc sumaDosArreglos.c -o sumaDosArreglos.out
\end{lstlisting}

Por ejemplo, la ejecución con siete procesos se realiza mediante:

\begin{lstlisting}
mpirun -np 7 ./sumaDosArreglos.out
\end{lstlisting}

El programa genera e imprime dos arreglos de 100 enteros, reconstruye el
arreglo de resultados y comprueba la suma elemento a elemento. Como los datos
son pseudoaleatorios, los valores mostrados pueden cambiar entre ejecuciones.

\section{Conclusión}

El programa implementa un esquema de paralelización coordinador--trabajadores
mediante MPI para calcular la suma elemento a elemento de dos arreglos. El
proceso cero genera los operandos, distribuye sus segmentos, recibe los
resultados y reconstruye el arreglo final; mientras tanto, los trabajadores
calculan las sumas de sus propias secciones.

La implementación permite comprender el uso de \texttt{MPI\_Send()} y
\texttt{MPI\_Recv()} en el intercambio de arreglos, así como la importancia
de las etiquetas para distinguir los mensajes. También muestra cómo manejar
una distribución no exacta cuando el número de elementos no es divisible
entre la cantidad de trabajadores.

La verificación final asegura que, durante la ejecución comprobada, cada
posición del arreglo resultante coincida con la suma de sus dos operandos.
Con ello se demuestra la correcta aplicación de la descomposición de datos
y la recolección de resultados en este ejercicio de programación paralela.

\section{Capturas de comprobación}

\captura{codigoFuente.png}{
    Código fuente del programa de suma de dos arreglos en el editor NeoVim.
}	

\captura{terminal.png}{
    Compilación y ejecución del programa con distintos valores de \texttt{N}
    y diferentes cantidades de procesos, incluyendo la verificación del resultado.
}

\end{document}